\chapter{導入}
\label{ch:introduction}

Givens–Shortt の論文
\emph{A class of Wasserstein metrics for probability distributions}
（Michigan Math.\ J.\ \textbf{31} (1984), 231–240）の
Proposition 2 を目標とする：
\[
  1\leq p\leq\infty
  \quad\Longrightarrow\quad
  \Wass_p\text{ は }\Pp{p}(X)\text{ 上の距離である．}
\]

原論文と同じく $(X,d)$ は Polish 空間，
$\Prob(X)$ はその Borel 確率測度全体とする．
Polish 空間上の Borel 確率測度は自動的に Radon 測度なので，
Radon 性を別に仮定・準備する必要はない．
